\begin{table}[htbp]
\centering
\small
\begin{tabular}{rrrrrrrr}
\toprule
rule & $N$ & sectors & $D_{\max}$ & $D_{\max}/2^N$ & mean & median & singletons \\
\midrule
156 & 13 & 610 & 48 & 5.86e-03 & 13.4 & 18 & 8 \\
156 & 17 & 4\,181 & 144 & 1.10e-03 & 31.3 & 45 & 10 \\
156 & 21 & 28\,657 & 432 & 2.06e-04 & 73.2 & 108 & 12 \\
108 & 13 & 1\,897 & 233 & 2.84e-02 & 4.3 & 8 & 714 \\
108 & 17 & 17\,991 & 1\,597 & 1.22e-02 & 7.3 & 21 & 4\,895 \\
108 & 21 & 170\,625 & 10\,946 & 5.22e-03 & 12.3 & 42 & 33\,552 \\
\bottomrule
\end{tabular}
\caption{The numbers behind Figure~\ref{fig:r19hist} (\rt{obc0}). \emph{mean} is $2^N/n_{\wcc}$; \emph{median} is the size of the sector a uniformly random computational-basis state lands in, i.e.\ where the cumulative curve crosses $1/2$. Both are far below $D_{\max}$, which is why the largest sector is a poor description of the typical one.}
\label{tab:r19sizes}
\end{table}
